\scanpage{43}
\mspage{17}
% archive foliation: 45
En fait, les anti-involutions de $\mathbb C$ ont toutes des pts fixes, et ce sont les symétries orth. p.r. à des droites réelles (pas néc. passant par l'origine) dans $\mathbb C$. La surface à bord conforme correspondante est isom. à un \uuline{demi-plan} (avec bord). On note que tte métrique \inserted{sur $\mathcal X$,} quasi-invariante par autom. conformes, est invariante par autom. d'ordre fini, en particulier par $\sigma$. Donc il y a sur $X$ une classe \inserted{can.} de métriques riemanniennes, mod. cte.

\uline{NB} \deleted{\unsure{Le}} Le modèle de $X$ envisagé est le demi-plan supérieur \deleted{$\unsure{\mathbb C}$}
\[
\overline{\mathbb H}=\{z\in\mathbb C\mid\mathrm{Im}\,z\geq0\}.
\]
Le groupe des autom. de $\overline H$, isom. au groupe des autom. de la droite affine $\mathbb C$ qui commutent à $\sigma:z\longmapsto\overline z$, est \uline{le groupe affine réel}
